\title{浏览器端开源 CAD 软件开发}
\author{杨岢瑞}
\date{Mar 31, 2026}
\maketitle
计算机辅助设计（CAD）软件的发展历程堪称工程技术领域的传奇。从上世纪 60 年代的线描仪时代，到 80 年代 AutoCAD 等桌面软件的霸主地位，再到如今的云端协作与 Web 化转型，CAD 始终是制造业与建筑业的核心生产力工具。传统 CAD 软件如 AutoCAD 和 SolidWorks 虽然功能强大，但其安装过程繁琐、对硬件要求极高，且跨平台兼容性差，导致中小企业与个人设计师望而却步。这些局限性在移动互联网时代愈发凸显，用户迫切需要一种无需安装、随时随地可用的解决方案。\par
浏览器端 CAD 的兴起得益于 Web 技术的革命性进步。HTML5 Canvas 提供了高效的 2D 绘图能力，WebGL 则解锁了浏览器内的 3D 硬件加速渲染，而 WebAssembly（WASM）更是让复杂几何计算在浏览器中成为可能。这些技术合力，使得零安装、跨 PC/移动平台、实时多人协作以及无缝集成成为现实。以往需要数 GB 安装包的 CAD 功能，如今可在任意现代浏览器中流畅运行。市场需求同样强劲：中小企业希望降低软件采购成本，独立设计师追求移动端灵感捕捉，学生与教育机构则需要免费易用的教学工具。数据显示，全球 CAD 市场规模已超 100 亿美元，而 Web CAD 细分领域年增长率超过 30\%{}。\par
开源在这一领域的价值尤为突出。传统开源 CAD 如 FreeCAD（参数化 3D 建模）、LibreCAD（2D 绘图）和 OpenSCAD（脚本式建模）虽优秀，但仍局限于桌面环境。浏览器端开源 CAD 则开启了新纪元：社区驱动的快速迭代、低门槛贡献，以及免费分发让全球开发者能共同完善内核。想象一下，一个基于 WebAssembly 的 OpenCascade 端口，能让浏览器直接处理工业级 BREP 几何，这不仅是技术突破，更是开源精神的延续。\par
本文面向 Web 开发者、CAD 爱好者和开源贡献者，提供从概念到实战的完整开发指南。我们将剖析核心技术栈、审视现有项目、亲手构建一个功能齐全的浏览器 CAD，并分享开源最佳实践。无论你是想 fork 现有项目，还是从零打造属于自己的 CAD 工具，这篇文章都将为你铺平道路。通过 8000 余字的详尽剖析、实战代码与优化策略，你将掌握浏览器 CAD 的全栈开发之道。\par
\chapter{2. 浏览器端 CAD 的核心技术栈}
浏览器端 CAD 的灵魂在于高效渲染与精确计算的双轮驱动。渲染引擎首推 WebGL，它利用 GPU 进行硬件加速 3D 渲染，能轻松处理数百万顶点的复杂模型。相比 2D Canvas 的软件渲染，WebGL 的帧率可提升数十倍，但学习曲线较陡。为此，Three.js 和 Babylon.js 等库应运而生，它们封装了底层 WebGL API，提供场景管理、材质系统和光照计算。举例来说，使用 Three.js 实现一个基础 3D 建模视图非常直观。以下代码创建一个旋转的立方体场景：\par
\begin{lstlisting}[language=javascript]
import * as THREE from 'three';
import { OrbitControls } from 'three/examples/jsm/controls/OrbitControls.js';

const scene = new THREE.Scene();
const camera = new THREE.PerspectiveCamera(75, window.innerWidth / window.innerHeight, 0.1, 1000);
const renderer = new THREE.WebGLRenderer();
renderer.setSize(window.innerWidth, window.innerHeight);
document.body.appendChild(renderer.domElement);

const geometry = new THREE.BoxGeometry();
const material = new THREE.MeshBasicMaterial({ color: 0x00ff00 });
const cube = new THREE.Mesh(geometry, material);
scene.add(cube);

camera.position.z = 5;

const controls = new OrbitControls(camera, renderer.domElement);
controls.enableDamping = true;  // 惯性拖拽

function animate() {
  requestAnimationFrame(animate);
  cube.rotation.x += 0.01;
  cube.rotation.y += 0.01;
  controls.update();
  renderer.render(scene, camera);
}
animate();
\end{lstlisting}
这段代码首先导入 Three.js 核心模块和轨道控制器。创建场景、透视相机和 WebGL 渲染器，并设置画布尺寸。然后生成一个立方体几何体，使用基本材质着色并添加到场景。相机定位于 Z=5 处，确保模型可见。OrbitControls 提供鼠标拖拽旋转、缩放功能，enableDamping 属性添加平滑惯性效果。动画循环通过 requestAnimationFrame 持续旋转立方体并更新渲染。这是一个 CAD 视图的基础雏形，后续可扩展为模型编辑器。\par
计算引擎的瓶颈在于 JavaScript 的单线程与浮点精度不足，无法应对 CAD 的复杂几何运算，如布尔求交或 NURBS 曲面拟合。此时 WebAssembly 登场，它允许将 C++ 或 Rust 代码编译为浏览器可执行的二进制模块，性能接近原生应用。典型应用是移植 OpenCascade（OCCT），业界领先的开源 CAD 内核，支持 BREP（边界表示）建模。opencascade.js 便是其 WASM 端口，能在浏览器中执行 STEP 文件解析与实体运算。性能对比显示，WASM 版布尔运算速度是纯 JS 的 50 倍以上，这对于实时预览至关重要。\par
几何建模依赖专用库。OpenCascade 通过 WASM 提供工业级内核，包括精确的曲线曲面表示和拓扑操作。CGAL.js 专注计算几何，如凸包与三角剖分，但仍处实验阶段。Three-CSG 则用纯 JS 实现构造实体几何（CSG），适合快速原型。earcut 库专攻多边形三角剖分，高性能源于增量算法，避免了昂贵的 Delaunay 三角化。对于 UI 交互，React 或 Vue 结合 Ant Design 构建工具栏与属性面板。拖拽操作可借助 Konva.js 的事件系统，或自定义 mouse/touch 事件处理程序，确保触屏适配。Hammer.js 进一步简化多点触控手势，如捏合缩放。\par
数据格式支持是 CAD 的命脉。浏览器需处理 STEP/IGES（精确交换）、STL/OBJ（网格）和 SVG（2D）。Three.js 自带 loaders 模块，能异步加载这些格式，结合 FileSaver.js 实现导出。实际开发中，先验证文件 MIME 类型，再用 WASM 解析精确数据，最后转换为 WebGL 可渲染网格。这种技术栈组合，确保了浏览器 CAD 的高保真与高效。\par
\chapter{3. 现有浏览器端开源 CAD 项目分析}
浏览器端开源 CAD 项目已初具规模，但各有侧重。OpenJSCAD 以 5k+ GitHub Stars 领跑，支持脚本式参数化建模，技术栈融合 JS 和 WASM，社区活跃度高。CadHub 则提供桌面 Web 双端体验，集成 Three.js 和多种内核，活跃度中等。FreeCAD Web 是实验性移植，借力 WASM 运行 FreeCAD 内核，但活跃度较低。Sketchlet 聚焦 2D 矢量编辑，使用 SVG 与 Canvas，新兴但势头强劲。这些项目奠定了基础，却也暴露共性痛点。\par
以 OpenJSCAD 为例，其架构精妙：左侧在线代码编辑器实时编译用户脚本，右侧 Three.js 预览 3D 模型。源码亮点在于参数化 DSL，用户可编写如 function main() \{{} return cylinder(\{{}h: 20, r: 10\}{}); \}{} 的声明式代码，内核自动展开为 CSG 树。浏览器端通过 jscad-openjscad 渲染器处理，此 DSL 借鉴 OpenSCAD，易学且强大。但交互性是其短板：缺乏直观的拖拽编辑，只能靠代码迭代，限制了非程序员用户。\par
CadHub 的创新在于多内核支持，既兼容 OpenSCAD 脚本，也集成 OpenCascade WASM 端口。部署采用 Docker 容器化，确保一致性，用户上传 STEP 文件后，即可参数化修改并导出。架构上，它用 Electron 桥接桌面与 Web，Three.js 负责视图，WASM offload 计算。这种混合模式扩展性强，但浏览器纯 Web 版加载时间较长，受限于 WASM 模块体积。\par
这些项目的痛点总结为三点：性能瓶颈源于网格密集模型的渲染压力，内核移植难度高（如 FreeCAD 的完整约束求解未落地），文件兼容性差（STEP 解析常丢失拓扑信息）。此外，缺乏标准化 UI 与实时协作，进一步阻碍主流采用。开发者需从中汲取教训，在新项目中优先 WASM 多线程与渐进式加载。\par
\chapter{4. 从零开发浏览器端 CAD 软件实战指南}
实战从项目初始化开始。使用 Vite 构建 React + TypeScript 脚手架，能快速启动开发服务器并享受热重载。执行 npm create vite@latest cad-web -- --template react-ts，进入目录后安装核心依赖：npm i three @types/three opencascade.js mathjs fabricjs react-three-fiber。这些包覆盖渲染、几何与 UI 需求。Vite 的 ESM 打包确保 WASM 模块无缝集成。\par
首个核心模块是 3D 视图控制器。基于 Three.js OrbitControls，实现轨道导航与截面视图。扩展代码如下：\par
\begin{lstlisting}[language=tsx]
import React, { useRef, useEffect } from 'react';
import * as THREE from 'three';
import { OrbitControls } from 'three/examples/jsm/controls/OrbitControls';
import { Canvas, useThree } from '@react-three/fiber';

function Scene() {
  const { camera, gl } = useThree();
  const controlsRef = useRef<OrbitControls>();

  useEffect(() => {
    const controls = new OrbitControls(camera, gl.domElement);
    controlsRef.current = controls;
    controls.enableDamping = true;
    controls.dampingFactor = 0.05;
    // 截面视图：clipping planes
    camera.layers.enable(1);  // 启用层 1 用于截面
  }, [camera, gl]);

  return null;
}

export default function Viewport() {
  return (
    <Canvas camera={{ position: [0, 0, 10] }}>
      <Scene />
      <mesh>
        <boxGeometry args={[2, 2, 2]} />
        <meshStandardMaterial color="hotpink" />
      </mesh>
    </Canvas>
  );
}
\end{lstlisting}
这段 React 组件利用 react-three-fiber（R3F）声明式渲染 Three.js 场景。useThree 钩子访问相机与渲染器引用。useEffect 初始化 OrbitControls，dampingFactor 控制阻尼强度。新增 clipping planes 支持截面视图，通过 camera.layers 隔离渲染层，实现动态剖切。这为后续模型检查奠基，性能上 R3F 的 hooks 避免了手动循环。\par
第二个模块是 2D 草图编辑器。使用 Fabric.js 实现约束绘图，支持直线、圆弧与简单约束求解。约束如平行/垂直用向量点积判断：\par
\begin{lstlisting}[language=javascript]
import { fabric } from 'fabric';

const canvas = new fabric.Canvas('sketch-canvas', { width: 800, height: 600 });

function addLine(x1: number, y1: number, x2: number, y2: number) {
  const line = new fabric.Line([x1, y1, x2, y2], {
    stroke: 'black',
    strokeWidth: 2,
    selectable: true,
  });
  canvas.add(line);
}

// 约束求解：平行检查
function constrainParallel(line1: fabric.Line, line2: fabric.Line) {
  const vec1 = { x: line1.x2 - line1.x1, y: line1.y2 - line1.y1 };
  const vec2 = { x: line2.x2 - line2.x1, y: line2.y2 - line2.y1 };
  const dot = vec1.x * vec2.x + vec1.y * vec2.y;
  const mag1 = Math.sqrt(vec1.x**2 + vec1.y**2);
  const mag2 = Math.sqrt(vec2.x**2 + vec2.y**2);
  const cosTheta = dot / (mag1 * mag2);
  if (Math.abs(cosTheta - 1) < 0.01) {  // 余弦接近 1 为平行
    line2.set({ stroke: 'blue' });  // 视觉反馈
  }
  canvas.renderAll();
}
\end{lstlisting}
Fabric.js 的对象模型允许实时编辑路径。addLine 创建可选中线段，constrainParallel 计算方向向量夹角，若余弦绝对值接近 1，则标记为平行约束（蓝色高亮）。实际中，可集成完整求解器如 Sylvester 库处理过度约束系统。此模块输出 SVG 路径，供 3D 拉伸使用。\par
第三个模块聚焦几何建模，实现拉伸、旋转与布尔运算。Three-CSG 处理 CSG 操作：\par
\begin{lstlisting}[language=javascript]
import * as THREE from 'three';
import * as CSG from 'three-csg-ts';

function extrudeProfile(profile: THREE.Shape, height: number): THREE.Mesh {
  const geometry = new THREE.ExtrudeGeometry(profile, { depth: height });
  const material = new THREE.MeshStandardMaterial({ color: 0x44ff44 });
  return new THREE.Mesh(geometry, material);
}

async function booleanSubtract(a: THREE.Mesh, b: THREE.Mesh): Promise<THREE.Mesh> {
  const aGeom = a.geometry.clone();
  const bGeom = b.geometry.clone();
  const result = CSG.subtract(aGeom, bGeom);
  const material = new THREE.MeshStandardMaterial({ color: 0xff4444 });
  return new THREE.Mesh(result, material);
}
\end{lstlisting}
extrudeProfile 从 2D 轮廓拉伸实体，ExtrudeGeometry 参数 depth 控制高度。booleanSubtract 使用 three-csg-ts 库的 subtract 方法，克隆几何体避免修改原 mesh，返回差集结果。CSG 算法基于 BSP 树，浏览器端通过 Worker 异步执行，避免 UI 阻塞。\par
参数化系统用 math.js 解析表达式。用户输入如「高度 = 长度 * 1.5」，math.js 求值为约束传播：\par
\begin{lstlisting}[language=javascript]
import { create, all } from 'mathjs';

const { evaluate, parse } = create(all);

interface Param {
  name: string;
  expr: string;
  value: number;
}

const params: Param[] = [];

function updateParam(name: string, expr: string) {
  const node = parse(expr);
  const scope = params.reduce((acc, p) => ({ ...acc, [p.name]: p.value }), {});
  const value = evaluate(node, scope);
  const param = params.find(p => p.name === name) || { name, expr, value: 0 };
  param.expr = expr;
  param.value = value;
  // 触发几何重构
  rebuildModel();
}
\end{lstlisting}
math.js 的 parse 生成抽象语法树，evaluate 在作用域中求值，支持循环依赖检测。此系统动态更新模型尺寸，实现 FreeCAD 式参数化。\par
文件 IO 模块集成 opencascade.js 解析 STEP：\par
\begin{lstlisting}[language=javascript]
import { OC } from 'opencascade.js';

async function loadSTEP(file: File): Promise<THREE.Group> {
  const buffer = await file.arrayBuffer();
  const stepData = new OC.STEPControl_Reader_();
  stepData.ReadFile('NUL', new Uint8Array(buffer));  // OCCT API
  stepData.TransferRoots();
  const shape = stepData.OneShape();
  // 转换为 Three.js meshes（简化）
  const meshes = ocToThree(shape);
  return new THREE.Group().add(...meshes);
}
\end{lstlisting}
OCCT 的 STEPControl\_{}Reader 处理精确 BREP，TransferRoots 导入根实体，OneShape 合并为单一形状。自定义 ocToThree 函数遍历拓扑，生成 WebGL 网格。\par
性能优化至关重要。LOD 通过 THREE.LOD 对象动态切换细节级别：远距离用低聚模型。Worker 线程 offload CSG 计算：\par
\begin{lstlisting}[language=javascript]
const worker = new Worker('csg-worker.js');
worker.postMessage({ geomA: aGeom.toJSON(), geomB: bGeom.toJSON() });
worker.onmessage = (e) => {
  const resultGeom = THREE.BufferGeometry.fromJSON(e.data);
  scene.add(new THREE.Mesh(resultGeom, material));
};
\end{lstlisting}
内存管理调用 geometry.dispose() 和 material.dispose() 释放 GPU 资源。测试用 Vitest 覆盖单元场景，Lighthouse 审计性能指标。\par
完整 Demo 仓库：https://github.com/yourname/cad-web-template，提供一键部署模板。\par
\chapter{5. 开源最佳实践与社区建设}
开源成功离不开许可策略。MIT 许可推荐指数最高，其宽松条款便于商业集成，仅要求保留版权声明。相比 LGPL（链接时需开源），MIT 更吸引企业用户，推动生态扩张。\par
文档是项目的门面。README 需包含安装、快速上手与架构图，在线 Demo 用 Vercel 一键部署，确保浏览器即用。API 文档借助 Storybook，交互式展示组件如视图控制器。\par
贡献指南标准化流程：Issue 模板分类 Bug/Feature，PR 要求关联 Issue 并通过 CI 检查。代码规范用 ESLint（airbnb 规则）与 Prettier 统一风格，husky 钩子强制提交前格式化。\par
发布路径多样：npm 发布核心库如 cad-core，Web 打包为静态站点。PWA 支持通过 manifest.json 与 Service Worker，实现离线编辑。\par
变现不违背开源精神：SaaS 提供云存储与协作，插件市场售卖专业模块，GitHub Sponsors 获社区赞助。\par
浏览器 CAD 仍面临挑战。性能需 WASM 多线程（SharedArrayBuffer）与 GPU 计算（WebGPU）。精度问题源于 JS 浮点误差，使用 BigFloat 库或 OCCT 高精度模式缓解。兼容性通过 Polyfill（如 core-js）桥接旧浏览器。安全上，Sandbox 用户脚本执行，避免 eval 注入。\par
未来趋势激动人心。AI 辅助设计集成 TensorFlow.js，实现草图自动识别：$\mathbf{y} = f(\mathbf{x}; \theta)$ 神经网络预测几何参数。WebXR 开启 AR/VR 建模，Yjs + WebRTC 实现 OT（操作转换）实时协作。glTF 2.0 标准化 CAD 传输，压缩率达 90\%{}。\par
行动起来！fork 本文 Demo，贡献你的几何模块。资源包括 OpenCascade 文档、Three.js 示例与 CAD 论坛。\par
浏览器端开源 CAD 将重塑设计范式，开源社区是先锋力量。加入我们，共筑 Web CAD 时代！\par
\chapter{附录}
\textbf{A. 完整技术栈清单}\par
three@0.158.0, opencascade.js@0.18.1, mathjs@12.4.2 等。安装：npm i three@0.158.0 等。\par
\textbf{B. 参考文献}\par
《WebGL Programming Guide》，OpenJSCAD GitHub，OCCT 文档。\par
\textbf{C. 术语表}\par
BREP：边界表示几何。NURBS：非均匀有理 B 样条。约束求解：几何关系闭合系统。\par
\textbf{D. 更新日志}\par
v1.0：初版。计划季度更新 WebGPU 支持。\par
